
We end by giving an application, to stable homotopy theory, of Burklund's Theorem \ref{thm:Burklund}. Fix a prime number $p$.

\begin{dfn}
For each integer $n \ge 1$, Serre proved \cite{SerreFinite} that the $p$-local $n$th stable stem 
$\pi_n(\mathbb{S}_{(p)})$
is a finite $p$-group.  The \emph{exponent} of the $n$th stable stem, denoted here by
$$\mathrm{exp}(\pi_n(\mathbb{S}_{(p)})),$$
is the smallest integer $a$ such that all elements of $\pi_n(\mathbb{S}_{(p)})$ are $p^a$-torsion. 
\end{dfn}

Upper bounds for the exponent have been considered in several papers \cite{Ad:SHT, liul, Arlettaz, Gonz00, MathewExponent}.  For $p>3$, the best prior bounds are due to Gonz\'alez \cite[Corollary 4.1.4]{Gonz00}. As in Gonz\'alez's work, our bounds on the exponent are deduced from an upper bound on $\Gamma_p$.
% who proves Adams spectral sequence vanishing lines to conclude his results.  Since Burklund obtains better asymptotic vanishing lines in Appendix A, he improves Gonz\'alez's exponent results:

\begin{thm}[Burklund] \label{thm:torsion-bound}
There is an inequality
$$\mathrm{exp}(\pi_n(\mathbb{S}_{(p)})) \le \frac{(2p-1)n}{(2p-2)(2p^2-2)} + o(n),$$
where $o(n)$ is the sublinear error term appearing in the statement of Theorem \ref{thm:Burklund}(1).
\end{thm}

\begin{proof}
At odd primes, Adams showed that the image of $J$ is a direct summand of $\pi_n(\mathbb{S}_{(p)})$ \cite{Ad:SHT, AdamsJIV}.  At the prime $2$, Adams and Quillen proved that the subgroup generated by the image of $J$ and the $\mu$-family is a direct summand \cite{QuillenAdamsConj}.  These papers also calculate the order of the image of $J$, from which it follows that the exponents of these summands grow logarithmically in $n$.

Suppose now that $x$ is an element of the complementary summand of $\pi_n(\mathbb{S}_{(p)})$, and let $\Gamma_p(n)$ denote the function from the statement of Theorem \ref{thm:Burklund}.  Since multiplication by $p$ raises $\mathrm{H}\mathbb{F}_p$-Adams filtration by at least $1$, Theorem \ref{thm:Burklund} implies that $p^{\Gamma_p(n)}x$ is either in the image of $J$, or, if $p=2$, in the subgroup generated by the image of $J$ and the $\mu$-family.  Since we assumed that $x$ is in the complementary summand, it follows that $p^{\Gamma_p(n)}x=0$.
\end{proof}

\begin{rmk}
  Like previous bounds on the torsion exponent of the stable stems,
  \Cref{thm:torsion-bound} is a linear bound.
  By contrast it is expected that $\mathrm{exp}(\pi_n(\mathbb{S}_{(p)}))$ grows sublinearly in $n$,
  though it remains unclear what the specific asymptotics of this function should be.
  The best known lower bound is logarithmic and comes from the image of $J$.
\end{rmk}

% An elaboration of the above result, with a more explicit discussion of the error term, may be found in Appendix \ref{sec:Appendix}.

%%% Local Variables:
%%% mode: latex
%%% TeX-master: "Main"
%%% eval: (visual-line-mode t)
%%% eval: (auto-fill-mode 0)
%%% End
